\section{Model and NTK framework}

\vspace{1em}
\subsection{RF-LR architecture (simplified)}
Let \( \bm{h}^{(0)}(\bm{x})=\bm{x}\in\mathbb{R}^{d_0} \). For \( \ell=1,\dots,L \),
\begin{equation}
\label{eq:RF-LR}
\bm{h}^{(\ell)}(\bm{x})
= \frac{1}{\sqrt{n_\ell}} \sum_{j=1}^{n_\ell} \bm{A}^{(\ell)}_j\,
\ReLU\!\Big(\bm{w}^{(\ell)\top}_j\,\bm{h}^{(\ell-1)}(\bm{x})\Big).
\end{equation}
\textbf{Training policy:} \( \bm{A}^{(\ell)} \) are trained; \( \bm{w}^{(\ell)} \) are frozen random draws (i.i.d. Gaussian). The \( 1/\sqrt{n_\ell} \) scaling (NTK parameterization) ensures a well-defined infinite-width limit. Low-rank bottlenecks (intermediate dimension \( r\ll n_\ell \)) compress correlations, reduce parameter budget from \( O(n^2) \) to \( O(rn) \), and improve conditioning.

\textbf{Remark (NTK magnitude with partial training):} By training only output weights \( \bm{A}^{(\ell)} \) while freezing input weights \( \bm{w}^{(\ell)} \), the effective NTK magnitude is approximately \emph{half} of what it would be if all parameters were trained. This is because the NTK decomposes as a sum over trainable parameter contributions; freezing \( \bm{w}^{(\ell)} \) removes their gradient contribution. Consequently, in the kernel regression regime, the effective learning rate is doubled (or equivalently, the kernel eigenvalues are halved), which alters convergence speed but not the final solution in the RKHS. This factor-of-2 effect on kernel magnitude should be accounted for when comparing RF-LR training dynamics to standard fully-trainable networks. (TO BE FORMALIZED: explicit computation showing \( \Kop_{\text{RF-LR}} \approx \tfrac{1}{2} \Kop_{\text{full}} \) for comparable architectures.)

\paragraph{Why this simplification?}
Removing bias scaling \( \beta \) and output scaling \( (\sigma_A, \sigma_c) \) yields cleaner expressions without loss of generality: (i) bias can be absorbed via data augmentation \( x \mapsto [x, 1] \), (ii) output scaling can be absorbed into learning rate, (iii) the mean-field and NTK structures remain qualitatively identical, and (iv) the Fourier-space interpretation and stepwise loss phenomenon depend on kernel geometry, not absolute scales.

\vspace{1em}
\subsection{Base NNGP and derivative kernels}
In the sequential infinite-width limit, the layer-\( \ell \) base kernel is
\begin{equation}
\Sigma^{(\ell)}(x,x')
= \mathbb{E}_{\bm{w}}\Big[\ReLU\!\big(\bm{w}^\top \bm{h}^{(\ell-1)}(x)\big)\,\ReLU\!\big(\bm{w}^\top \bm{h}^{(\ell-1)}(x')\big)\Big],
\end{equation}
and the derivative kernel is
\begin{equation}
\dot{\Sigma}^{(\ell)}(x,x')
= \mathbb{E}_{\bm{w}}\Big[\dot{\ReLU}\!\big(\bm{w}^\top \bm{h}^{(\ell-1)}(x)\big)\,\dot{\ReLU}\!\big(\bm{w}^\top \bm{h}^{(\ell-1)}(x')\big)\;\bm{w}^\top\bm{w}\Big].
\end{equation}
For \( \ell=1 \) and centered Gaussian \( \bm{w} \), \( \Sigma^{(1)} \) is rotation-invariant with standard closed forms in terms of input cosine similarity \( \rho = \langle x,x'\rangle/(\|x\|\|x'\|) \):
\[
\Sigma^{(1)}(x,x') = \frac{\|x\|\|x'\|}{2\pi}\big((\pi-\theta)\cos\theta + \sin\theta\big), \quad \theta = \arccos(\rho).
\]

\vspace{1em}
\subsection{NTK recursion}
\begin{theorem}[Infinite-width NTK recursion]
\label{thm:ntk_recursion}
Under NTK parameterization \( 1/\sqrt{n_\ell} \) and a sequential limit \( n_1\to\infty,\dots,n_L\to\infty \) with \( \bm{w}^{(\ell)} \sim \mathcal{N}(0, \mathbf{I}_{d_{\ell-1}}) \) i.i.d. frozen after initialization, the layer-\( \ell \) NTK satisfies
\begin{align}
\Kop^{(0)}(x,x') \;&=\; 0, \\
\Kop^{(\ell)}(x,x') \;&=\; \Kop^{(\ell-1)}(x,x')\cdot\dot{\Sigma}^{(\ell)}(x,x')\;+\;\Sigma^{(\ell)}(x,x'),\quad \ell\ge1,
\end{align}
where
\begin{align}
\Sigma^{(\ell)}(x,x') &= \mathbb{E}_{\bm{w}}\Big[\ReLU\!\big(\bm{w}^\top \bm{h}^{(\ell-1)}(x)\big)\,\ReLU\!\big(\bm{w}^\top \bm{h}^{(\ell-1)}(x')\big)\Big], \\
\dot{\Sigma}^{(\ell)}(x,x') &= \mathbb{E}_{\bm{w}}\Big[\dot{\ReLU}\!\big(\bm{w}^\top \bm{h}^{(\ell-1)}(x)\big)\,\dot{\ReLU}\!\big(\bm{w}^\top \bm{h}^{(\ell-1)}(x')\big)\;\bm{w}^\top\bm{w}\Big].
\end{align}
For \( \ell=1 \), \( \Sigma^{(1)} \) is deterministic; for \( \ell>1 \), \( \Sigma^{(\ell)} \) and \( \dot{\Sigma}^{(\ell)} \) are random fields whose fluctuations shrink with width and with the rank \( r \) of bottlenecks.

Furthermore, the gradient flow dynamics in the NTK regime satisfy
\[
\frac{d}{dt}f_t(x) = -\int \Kop^{(L)}(x,x') \big(f_t(x') - y(x')\big) d\mu(x'),
\]
where \( \mu \) is the data distribution and \( y \) is the target function.
\end{theorem}

Proof: See Appendix A.1.

\paragraph{Interpretation.} The additive term \( \Sigma^{(\ell)} \) encodes fresh basis contribution at layer \( \ell \), while the multiplicative coupling \( \Kop^{(\ell-1)}\cdot\dot{\Sigma}^{(\ell)} \) transports geometry through layers via the derivative kernel.

\medskip
\noindent\textbf{Corollary (Two-layer NTK).}
Consider \( L=2 \). Then
\begin{align*}
\Kop^{(1)}(x,x') \;&=\; \Sigma^{(1)}(x,x'),\\
\Kop^{(2)}(x,x') \;&=\; \Kop^{(1)}(x,x')\cdot\dot{\Sigma}^{(2)}(x,x') \;+\; \Sigma^{(2)}(x,x'),
\end{align*}
where \( \Sigma^{(1)}(x,x')=\mathbb{E}_{\bm{w}}\big[\ReLU(\bm{w}^\top x)\ReLU(\bm{w}^\top x')\big] \) and
\[
\dot{\Sigma}^{(2)}(x,x')=\mathbb{E}_{\bm{w}}\Big[\dot{\ReLU}\!\big(\bm{w}^\top \bm{h}^{(1)}(x)\big)\,\dot{\ReLU}\!\big(\bm{w}^\top \bm{h}^{(1)}(x')\big)\,\bm{w}^\top\bm{w}\Big],\quad
\Sigma^{(2)}(x,x')=\mathbb{E}_{\bm{w}}\Big[\ReLU(\cdot)\ReLU(\cdot)\Big].
\]
Proof: See Appendix A.2.

\medskip
\noindent\textbf{Corollary (General L-layer NTK).}
\label{cor:general_L_layer_ntk}
For any depth \( L \geq 1 \), the NTK admits the explicit form:
\begin{equation}
\label{eq:ntk_explicit_form}
\Kop^{(L)}(x,x') = \sum_{\ell=1}^{L} \Sigma^{(\ell)}(x,x') \prod_{k=\ell+1}^{L} \dot{\Sigma}^{(k)}(x,x'),
\end{equation}
where we use the convention that an empty product equals 1 (so the \( \ell=L \) term is just \( \Sigma^{(L)}(x,x') \)).

Equivalently, expanding the recursion yields:
\begin{align}
\Kop^{(L)}(x,x') &= \Sigma^{(L)}(x,x') + \Sigma^{(L-1)}(x,x')\cdot\dot{\Sigma}^{(L)}(x,x') \nonumber\\
&\quad + \Sigma^{(L-2)}(x,x')\cdot\dot{\Sigma}^{(L-1)}(x,x')\cdot\dot{\Sigma}^{(L)}(x,x') \nonumber\\
&\quad + \cdots + \Sigma^{(1)}(x,x')\cdot\prod_{k=2}^{L}\dot{\Sigma}^{(k)}(x,x').
\end{align}

\paragraph{Interpretation.}
Each term in the sum corresponds to a path through the network:
\begin{itemize}
    \item The term \( \Sigma^{(\ell)} \prod_{k=\ell+1}^L \dot{\Sigma}^{(k)} \) represents learning at layer \( \ell \) (fresh basis contribution \( \Sigma^{(\ell)} \)) propagated through layers \( \ell+1, \dots, L \) (derivative kernel transport \( \dot{\Sigma}^{(k)} \)).
    \item The deepest term \( \Sigma^{(L)} \) is learning at the final layer without transport.
    \item The shallowest term \( \Sigma^{(1)} \prod_{k=2}^L \dot{\Sigma}^{(k)} \) is learning at the first layer transported through all subsequent layers.
\end{itemize}

\paragraph{Depth dependence and scaling.}
The NTK magnitude scales roughly as:
\[
\Kop^{(L)}(x,x') \sim L \cdot \bar{\Sigma} \cdot (\bar{\dot{\Sigma}})^{L-1},
\]
where \( \bar{\Sigma} \) and \( \bar{\dot{\Sigma}} \) are typical magnitudes of the base and derivative kernels. For ReLU networks with normalized inputs:
\begin{itemize}
    \item \( \bar{\Sigma} = O(1) \) (bounded by input norm).
    \item \( \bar{\dot{\Sigma}} = O(1) \) but typically \( < 1 \), causing exponential decay with depth.
    \item The factor \( L \) provides polynomial growth, but \( (\bar{\dot{\Sigma}})^{L-1} \) causes exponential decay.
\end{itemize}
This explains the vanishing gradient problem in deep untrained networks and motivates careful initialization schemes to maintain \( \bar{\dot{\Sigma}} \approx 1 \).

Proof: See Appendix A.3.
